\section{Quantum shuf\/f\/le algebras}\label{sec:shuffle}

The quantized coordinate ring $\mathcal{A}_v[N]$ admits an embedding into a~quantum shuf\/f\/le algebra associated to the
$I\times I$ symmetrizable Cartan matrix~$A$ with symmetrizers $\mathbf{d}$.
In this section we introduce our choice of quantum shuf\/f\/le algebra and prove some basic properties.
Better proofs of many of these results can be found in~\cite{leclerc}.
We refer the reader also to \cite{Rosso} for a more general discussion of quantum shuf\/f\/le algebras.


Let $\mathbf{W}=\bigsqcup\limits_{m\ge0} I^m$ denote the set of all words in the alphabet~$I$.
For a~word $\mathbf{i}=(i_1,\ldots,i_m)$ and a~sequence of nonnegative integers $\mathbf{c}\in\mathbb{Z}_{\ge0}^m$ we
will make use of the following simplifying notation:
\begin{gather*}
\mathbf{i}^\mathbf{c}=(i_1^{c_1},\ldots,i_m^{c_m})=(\underbrace{i_1,\ldots,
i_1}_{c_1},\ldots,\underbrace{i_m,\ldots,i_m}_{c_m}),
\end{gather*}
where $i_k$ appears $c_k$ times ($1\le k\le m$).
For $\mathbf{j}=(j_1,\ldots,j_r)\in\mathbf{W}$ we will write $|\mathbf{j}|=\alpha_{j_1}+\dots+\alpha_{j_r}$ for the
\emph{degree} of the word $\mathbf{j}$.
Denote by $\mathbf{W}^\nu$ the set of all words of degree $\nu\in\mathcal{Q}$.

Def\/ine the quantum shuf\/f\/le algebra $g^*$ as the $\mathbb{Z}[v^{\pm\frac{1}{2}}]$-algebra with basis $\mathbf{W}$ and
multiplication given~by
\begin{gather}
\label{eq:shuffle_product}
(j_1,\ldots,j_r)\circ(j_{r+1},\ldots,j_{r+s})=\sum\limits_{\sigma\in\Sigma_{r,s}}v^{\zeta(\sigma)}\mathbf{j}_\sigma,
\end{gather}
where $\mathbf{j}_\sigma=(j_{\sigma_1},\ldots,j_{\sigma_{r+s}})$ and $\zeta(\sigma)=\zeta_{r,s}(\sigma)$ measures both
inversions and non-inversions occurring in~$\sigma$:
\begin{gather*}
\zeta(\sigma)=\frac{1}{2}\sum\limits_{\substack{1\le k\le r
\\
r+1\le \ell\le r+s
\\
\sigma^{-1}_\ell<\sigma^{-1}_k}}(\alpha_{j_k},\alpha_{j_\ell})-\frac{1}{2}\sum\limits_{\substack{1\le k\le r
\\
r+1\le \ell\le r+s
\\
\sigma^{-1}_k<\sigma^{-1}_\ell}}(\alpha_{j_k},\alpha_{j_\ell}).
\end{gather*}

\begin{Remark}
Kleshchev--Ram~\cite{kleshchev-ram1} and Leclerc~\cite{leclerc} use opposite multiplications for their quantum shuf\/f\/le
algebras measuring only inversions and only non-inversions respectively.
The conventions of~\cite{kleshchev-ram1} are claimed to be more convenient for the relationship with the representation
theory of quiver Hecke algebras.

We hope that our choice of multiplication on $g^*$ will prove equally useful in this regard since this is our preferred
quantum shuf\/f\/le multiplication which will f\/igure prominently in the results presented in this note.
\end{Remark}

The following proposition can be obtained from the results of~\cite{leclerc} by rescaling the basis, however for
completeness we provide a~proof.

\begin{Proposition}%\label{prop:shuffle_associativity}
The shuffle product~\eqref{eq:shuffle_product} defines an associative algebra structure on $g^*$.
\end{Proposition}

\begin{proof}
It will be convenient to introduce the subgroup $\Sigma_{r,s,t}$ of $\Sigma_{r+s+t}$ consisting of all shuf\/f\/les of the
sequences $(1,\ldots,r)$, $(r+1,\ldots,r+s)$, and $(r+s+1,\ldots,r+s+t)$, that is
\begin{gather*}
\Sigma_{r,s,t}=\big\{\sigma\in\Sigma_{r+s+t}: \sigma^{-1}_1<\dots<\sigma^{-1}_r, \sigma^{-1}_{r+1}<\dots<\sigma^{-1}_{r+s},
\sigma^{-1}_{r+s+1}<\dots<\sigma^{-1}_{r+s+t}\big\}.
\end{gather*}
There are natural embeddings $\Sigma_{r,s}\hookrightarrow\Sigma_{r,s,t}$ written $\sigma\mapsto\widetilde\sigma$ and
$\Sigma_{s,t}\hookrightarrow\Sigma_{r,s,t}$ written $\sigma\mapsto\widetilde\sigma'$, where
\begin{gather*}
\widetilde\sigma_k=\sigma_k
\quad
\text{for}
\quad
1\le k\le r+s
\qquad
\text{and}
\qquad
\widetilde\sigma_{r+s+k}=r+s+k
\quad
\text{for}
\quad
1\le k\le t,
\\
\widetilde\sigma'_k=k
\quad
\text{for}
\quad
1\le k\le r
\qquad
\text{and}
\qquad
\widetilde\sigma'_{r+k}=r+\sigma_k
\quad
\text{for}
\quad
1\le k\le s+t.
\end{gather*}
Notice that every element of $\Sigma_{r,s,t}$ has a~unique presentation as $\widetilde\sigma\cdot\tau$ with
$\sigma\in\Sigma_{r,s}$ and $\tau\in\Sigma_{r+s,t}$ and a~unique presentation as $\widetilde\sigma'\cdot\rho$ with
$\sigma\in\Sigma_{s,t}$ and $\rho\in\Sigma_{r,s+t}$.
For $\sigma\in\Sigma_{r,s,t}$ def\/ine
\begin{gather*}
\widetilde\zeta(\sigma)=\frac{1}{2}\sum\limits_{\substack{1\le k\le r
\\
r+1\le\ell\le r+s
\\
\sigma_\ell^{-1}<\sigma_k^{-1}}}(\alpha_{j_k},\alpha_{j_\ell}) +\frac{1}{2}\sum\limits_{\substack{1\le k\le r
\\
r+s+1\le\ell\le r+s+t
\\
\sigma_\ell^{-1}<\sigma_k^{-1}}}(\alpha_{j_k},\alpha_{j_\ell}) +\frac{1}{2}\sum\limits_{\substack{r+1\le k\le r+s
\\
r+s+1\le\ell\le r+s+t
\\
\sigma_\ell^{-1}<\sigma_k^{-1}}}(\alpha_{j_k},\alpha_{j_\ell})
\\
\phantom{\widetilde\zeta(\sigma)=}{}
 -\frac{1}{2}\sum\limits_{\substack{1\le k\le r
\\
r+1\le\ell\le r+s
\\
\sigma_k^{-1}<\sigma_\ell^{-1}}}(\alpha_{j_k},\alpha_{j_\ell}) -\frac{1}{2}\sum\limits_{\substack{1\le k\le r
\\
r+s+1\le\ell\le r+s+t
\\
\sigma_k^{-1}<\sigma_\ell^{-1}}}(\alpha_{j_k},\alpha_{j_\ell}) -\frac{1}{2}\sum\limits_{\substack{r+1\le k\le r+s
\\
r+s+1\le\ell\le r+s+t
\\
\sigma_k^{-1}<\sigma_\ell^{-1}}}(\alpha_{j_k},\alpha_{j_\ell}).
\end{gather*}
Then one immediately checks that:
\begin{gather*}
\widetilde\zeta(\widetilde\sigma\cdot\tau)=\zeta(\sigma)+\zeta(\tau)
\qquad
\text{for any $\sigma\in\Sigma_{r,s}$ and $\tau\in\Sigma_{r+s,t}$,}
\\
\widetilde\zeta(\widetilde\sigma'\cdot\rho)=\zeta(\sigma)+\zeta(\rho)
\qquad
\text{for any $\sigma\in\Sigma_{s,t}$ and $\rho\in\Sigma_{r,s+t}$.}
\end{gather*}
We are now ready to prove the proposition.
Indeed, on the one hand we have
\begin{gather*}
\big((j_1,\ldots,j_r)\circ(j_{r+1},\ldots,j_{r+s})\big)\circ(j_{r+s+1},\ldots,j_{r+s+t})
\\
\qquad =\sum\limits_{\sigma\in\Sigma_{r,s}}v^{\zeta(\sigma)}\mathbf{j}_{\sigma}\circ(j_{r+s+1},\ldots,j_{r+s+t})
 =\sum\limits_{\substack{\sigma\in\Sigma_{r,s}\\ \tau\in\Sigma_{r+s,t}}}v^{\zeta(\sigma)+\zeta(\tau)}\mathbf{j}_{\widetilde\sigma\cdot\tau}
\end{gather*}
and on the other hand we have
\begin{gather*}
(j_1,\ldots,j_r)\circ\big((j_{r+1},\ldots,j_{r+s})\circ(j_{r+s+1},\ldots,j_{r+s+t})\big)
\\
\qquad =\sum\limits_{\sigma\in\Sigma_{s,t}}v^{\zeta(\sigma)}(j_1,\ldots,j_r)\circ\mathbf{j}_\sigma
 =\sum\limits_{\substack{\sigma\in\Sigma_{s,t}\\ \rho\in\Sigma_{r,s+t}}}v^{\zeta(\sigma)+\zeta(\rho)}\mathbf{j}_{\widetilde\sigma'\cdot\rho}.
\end{gather*}
But both of these expressions are equal to
$\sum\limits_{\sigma\in\Sigma_{r,s,t}}v^{\widetilde\zeta(\sigma)}\mathbf{j}_\sigma$.
\end{proof}

We will also need a~comultiplication on $g^*$.
This is def\/ined on a~word $\mathbf{h}\in\mathbf{W}$~by
\begin{gather*}
\Delta(\mathbf{h})=\sum\limits_{\mathbf{j}_1,\mathbf{j}_2\in\mathbf{W}:(\mathbf{j}_1,\mathbf{j}_2)=\mathbf{h}}
v^{\frac{1}{2}(|\mathbf{j}_1|,|\mathbf{j}_2|)}\mathbf{j}_1\otimes\mathbf{j}_2.
\end{gather*}
In the following result we consider $g^*\otimes g^*$ with the twisted multiplication
\begin{gather*}
(\mathbf{j}_{1,1}\otimes\mathbf{j}_{1,2})\circ(\mathbf{j}_{2,1}\otimes\mathbf{j}_{2,2})
=v^{(|\mathbf{j}_{1,2}|,|\mathbf{j}_{2,1}|)}(\mathbf{j}_{1,1}\circ\mathbf{j}_{2,1})\otimes(\mathbf{j}_{1,2}\circ\mathbf{j}_{2,2})
\end{gather*}
for all words $\mathbf{j}_{1,1}$, $\mathbf{j}_{1,2}$, $\mathbf{j}_{2,1}$, $\mathbf{j}_{2,2}\in\mathbf{W}$. %$, $
\begin{Lemma}
The map $\Delta:g^*\to g^*\otimes g^*$ defines a~coassociative comultiplication on $g^*$.
Moreover, $g^*$ is a~bialgebra with counit $\epsilon:g^*\to\mathbb{C}$ given by $\varepsilon(\mathbf{j})=0$ whenever
$\mathbf{j}\ne()$ is not the empty word and $\varepsilon\big(()\big)=1$.
\end{Lemma}

\begin{proof}
We will check explicitly the coassociativity and the compatibility of multiplication and comultiplication on $g^*$, the
other axioms of a~bialgebra structure on $g^*$ are straightforward and left to the reader.

We start with the coassociativity of~$\Delta$.
For a~word $\mathbf{h}\in\mathbf{W}$ we have
\begin{gather*}
(\Delta\otimes\operatorname{Id})\Delta(\mathbf{h})
=\sum\limits_{\mathbf{j}_1,\mathbf{j}_2\in\mathbf{W}:(\mathbf{j}_1,\mathbf{j}_2)=\mathbf{h}}
v^{\frac{1}{2}(|\mathbf{j}_1|,|\mathbf{j}_2|)}\Delta(\mathbf{j}_1)
\otimes\mathbf{j}_2
\\
\phantom{(\Delta\otimes\operatorname{Id})\Delta(\mathbf{h})}
=\sum\limits_{\substack{\mathbf{j}_1,\mathbf{j}_2\in\mathbf{W}:\\(\mathbf{j}_1,\mathbf{j}_2)=\mathbf{h}}}
v^{\frac{1}{2}(|\mathbf{j}_1|,|\mathbf{j}_2|)}
\sum\limits_{\substack{\mathbf{j}_{1,1},\mathbf{j}_{1,2}\in\mathbf{W}:\\(\mathbf{j}_{1,1},\mathbf{j}_{1,2})=\mathbf{j}_1}}
v^{\frac{1}{2}(|\mathbf{j}_{1,1}|,|\mathbf{j}_{1,2}|)}(\mathbf{j}_{1,1}\otimes\mathbf{j}_{1,2})
\otimes\mathbf{j}_2
\\
\phantom{(\Delta\otimes\operatorname{Id})\Delta(\mathbf{h})}
=\sum\limits_{\mathbf{j}_{1,1},\mathbf{j}_{1,2},\mathbf{j}_2\in\mathbf{W}:(\mathbf{j}_{1,1},\mathbf{j}_{1,2},\mathbf{j}_2)=\mathbf{h}}
v^{\frac{1}{2}(|\mathbf{j}_{1,1}|+|\mathbf{j}_{1,2}|,|\mathbf{j}_2|)}
v^{\frac{1}{2}(|\mathbf{j}_{1,1}|,|\mathbf{j}_{1,2}|)}\mathbf{j}_{1,1}\otimes\mathbf{j}_{1,2}\otimes\mathbf{j}_2,
\\
(\operatorname{Id}\otimes\Delta)\Delta(\mathbf{h})
=\sum\limits_{\mathbf{j}_1,\mathbf{j}_2\in\mathbf{W}:(\mathbf{j}_1,\mathbf{j}_2)=\mathbf{h}}
v^{\frac{1}{2}(|\mathbf{j}_1|,|\mathbf{j}_2|)}\mathbf{j}_1\otimes\Delta(\mathbf{j}_2)
\\
\phantom{(\operatorname{Id}\otimes\Delta)\Delta(\mathbf{h})}
=\sum\limits_{\substack{\mathbf{j}_1,\mathbf{j}_2\in\mathbf{W}:\\(\mathbf{j}_1,\mathbf{j}_2)=\mathbf{h}}}
v^{\frac{1}{2}(|\mathbf{j}_1|,|\mathbf{j}_2|)}
\sum\limits_{\substack{\mathbf{j}_{2,1},\mathbf{j}_{2,2}\in\mathbf{W}:\\(\mathbf{j}_{2,1},\mathbf{j}_{2,2})=\mathbf{j}_2}}
v^{\frac{1}{2}(|\mathbf{j}_{2,1}|,|\mathbf{j}_{2,2}|)}\mathbf{j}_1\otimes(\mathbf{j}_{2,1}\otimes\mathbf{j}_{2,2})
\\
\phantom{(\operatorname{Id}\otimes\Delta)\Delta(\mathbf{h})}
=\sum\limits_{\mathbf{j}_1,\mathbf{j}_{2,1},\mathbf{j}_{2,2},\in\mathbf{W}:(\mathbf{j}_1,\mathbf{j}_{2,1},\mathbf{j}_{2,2})=\mathbf{h}}
v^{\frac{1}{2}(|\mathbf{j}_1|,|\mathbf{j}_{2,1}|+|\mathbf{j}_{2,2}|)}v^{\frac{1}{2}(|\mathbf{j}_{2,1}|,|\mathbf{j}_{2,2}|)}
\mathbf{j}_1\otimes\mathbf{j}_{2,1}\otimes\mathbf{j}_{2,2}.
\end{gather*}
But both of these are equal to
\begin{gather*}
\sum\limits_{\mathbf{j}_1,\mathbf{j}_2,\mathbf{j}_3,\in\mathbf{W}:(\mathbf{j}_1,\mathbf{j}_2,\mathbf{j}_3)=\mathbf{h}}
v^{\frac{1}{2}(|\mathbf{j}_1|,|\mathbf{j}_2|)}v^{\frac{1}{2}(|\mathbf{j}_1|,|\mathbf{j}_3|)}
v^{\frac{1}{2}(|\mathbf{j}_2|,|\mathbf{j}_3|)}\mathbf{j}_1\otimes\mathbf{j}_2\otimes\mathbf{j}_3,
\end{gather*}
this establishes coassociativity.

To see the compatibility of the shuf\/f\/le multiplication on $g^*$ and~$\Delta$ we consider two words
$\mathbf{j}_1,\mathbf{j}_2\in\mathbf{W}$.
On one hand we have
\begin{gather*}
\Delta(\mathbf{j}_1\circ\mathbf{j}_2)=\sum\limits_{\sigma\in\Sigma_{r,s}}v^{\zeta(\sigma)}\Delta(\mathbf{j}_\sigma)
=\sum\limits_{\sigma\in\Sigma_{r,s}}v^{\zeta(\sigma)}
\sum\limits_{\mathbf{j}_{\sigma,1},\mathbf{j}_{\sigma,2}\in\mathbf{W}:(\mathbf{j}_{\sigma,1},\mathbf{j}_{\sigma,2})=\mathbf{j}_\sigma}
v^{\frac{1}{2}(|\mathbf{j}_{\sigma,1}|,|\mathbf{j}_{\sigma,2}|)}\mathbf{j}_{\sigma,1}\otimes\mathbf{j}_{\sigma,2}.
\end{gather*}
On the other hand we have
\begin{gather*}
\Delta(\mathbf{j}_1)\circ\Delta(\mathbf{j}_2)
=\left(\sum\limits_{\substack{\mathbf{j}_{1,1},\mathbf{j}_{1,2}\in\mathbf{W}:\\(\mathbf{j}_{1,1},\mathbf{j}_{1,2})=\mathbf{j}_1}}
v^{\frac{1}{2}(|\mathbf{j}_{1,1}|,|\mathbf{j}_{1,2}|)}\mathbf{j}_{1,1}\otimes\mathbf{j}_{1,2}\right)\!
\circ\!\left(\sum\limits_{\substack{\mathbf{j}_{2,1},\mathbf{j}_{2,2}\in\mathbf{W}:\\(\mathbf{j}_{2,1},\mathbf{j}_{2,2})=\mathbf{j}_2}}
v^{\frac{1}{2}(|\mathbf{j}_{2,1}|,|\mathbf{j}_{2,2}|)}\mathbf{j}_{2,1}\otimes\mathbf{j}_{2,2}\right)
\\
\qquad
=\sum\limits_{\substack{\mathbf{j}_{1,1},\mathbf{j}_{1,2},\mathbf{j}_{2,1},\mathbf{j}_{2,2}\in\mathbf{W}:\\
(\mathbf{j}_{1,1},\mathbf{j}_{1,2})=\mathbf{j}_1,(\mathbf{j}_{2,1},\mathbf{j}_{2,2})=\mathbf{j}_2}}
v^{\frac{1}{2}(|\mathbf{j}_{1,1}|,|\mathbf{j}_{1,2}|)}v^{\frac{1}{2}(|\mathbf{j}_{2,1}|,|\mathbf{j}_{2,2}|)}
v^{(|\mathbf{j}_{1,2}|,|\mathbf{j}_{2,1}|)}(\mathbf{j}_{1,1}\circ\mathbf{j}_{2,1})\otimes(\mathbf{j}_{1,2}\circ\mathbf{j}_{2,2})
\\
\qquad
=\sum\limits_{\substack{\mathbf{j}_{1,1},\mathbf{j}_{1,2},\mathbf{j}_{2,1},\mathbf{j}_{2,2}\in\mathbf{W}:\\
(\mathbf{j}_{1,1},\mathbf{j}_{1,2})=\mathbf{j}_1,(\mathbf{j}_{2,1},\mathbf{j}_{2,2})=\mathbf{j}_2}}
\sum\limits_{\sigma_1\in\Sigma_{r_1,s_1}}
\sum\limits_{\sigma_2\in\Sigma_{r_2,s_2}}v^{\frac{1}{2}(|\mathbf{j}_{1,1}|,|\mathbf{j}_{1,2}|)}
v^{\frac{1}{2}(|\mathbf{j}_{2,1}|,|\mathbf{j}_{2,2}|)}v^{(|\mathbf{j}_{1,2}|,|\mathbf{j}_{2,1}|)}
\\
\qquad\phantom{=}  \times
v^{\zeta(\sigma_1)}v^{\zeta(\sigma_2)}\mathbf{j}_{\sigma_1,1}\otimes\mathbf{j}_{\sigma_2,2},
\end{gather*}
where we wrote $r_i$ and $s_i$ ($i=1,2$) for the lengths of $\mathbf{j}_{i,1}$ and $\mathbf{j}_{i,2}$ respectively.
There is an obvious bijection between pairs consisting of a~shuf\/f\/le $\sigma\in\Sigma_{r,s}$ together with a~partition
$(\mathbf{j}_{\sigma,1},\mathbf{j}_{\sigma,2})$ of~$\mathbf{j}_\sigma$ and quadruples consisting of partitions of~$\mathbf{j}_1$ and~$\mathbf{j}_2$ together with shuf\/f\/les $\sigma_1\in\Sigma_{r_1,s_1}$ and
$\sigma_2\in\Sigma_{r_2,s_2}$.
Indeed, reading of\/f the contributions in $\mathbf{j}_{\sigma,1}$ and $\mathbf{j}_{\sigma,2}$ coming from $\mathbf{j}_1$
and $\mathbf{j}_2$ produce the desired partitions while the relative positions of these contributions produce the
desired shuf\/f\/les in $\Sigma_{r_1,s_1}$ and $\Sigma_{r_2,s_2}$.

Notice that under this bijection we have
\begin{gather*}
\zeta(\sigma)=\zeta(\sigma_1)+\zeta(\sigma_2)+\frac{1}{2}(|\mathbf{j}_{1,2}|,|\mathbf{j}_{2,1}|)-\frac{1}{2}(|\mathbf{j}_{1,1}|,|\mathbf{j}_{2,2}|)
\end{gather*}
and
\begin{gather*}
(|\mathbf{j}_{\sigma,1}|,|\mathbf{j}_{\sigma,2}|)
=(|\mathbf{j}_{1,1}|,|\mathbf{j}_{1,2}|)+(|\mathbf{j}_{1,1}|,|\mathbf{j}_{2,2}|)
+(|\mathbf{j}_{2,1}|,|\mathbf{j}_{1,2}|)+(|\mathbf{j}_{2,1}|,|\mathbf{j}_{2,2}|).
\end{gather*}
An easy inspection now shows that $\Delta(\mathbf{j}_1\circ\mathbf{j}_2)=\Delta(\mathbf{j}_1)\circ\Delta(\mathbf{j}_2)$
as claimed.
\end{proof}
